% META: {"id":"TEXP-CAG-EX-001","chapitre":"TEXP-COMPLEXES-ALGEBRE-GEOMETRIE","type_objet":"exercice","capacites":["TEXP-CAG-C2"],"capacites_codes":["C2"],"methodes":[],"parcours":1,"competences":["calculer"],"duree_min":10,"mode_creation":"ex_nihilo","sources_inspiration":[],"fichier_tex":"chapitres/TEXP-COMPLEXES-ALGEBRE-GEOMETRIE/exercices/TEXP-CAG-EX-001.tex","corrige_tex":"chapitres/TEXP-COMPLEXES-ALGEBRE-GEOMETRIE/corriges/TEXP-CAG-CO-001.tex","status":"approved","origin":"generated","status_history":["generated","audited","accepted","approved"],"release_acceptance":"RELEASE_OWNER_BATCH_ACCEPTANCE","acceptance_closure_digest":"sha256:9b3ccf9a81c5520fbb7e03b7b2d3e7bf2057a02834b6d908bf3be5f49f3b553f","acceptance_receipt_digest":"sha256:4fad86f0282e0fa4e0419cca1ad6379ae7af5a280c04a4a90880f90a1c044242"}
% BEGIN-VERIFY
% from sympy import I, symbols, Eq, solve
% z = symbols('z')
% sol = solve(Eq((1-I)*z, 4+2*I), z)
% assert sol == [1 + 3*I]
% END-VERIFY
\begin{exercice}{TEXP-CAG-EX-001}{1}{10}
Résoudre l'équation $(1-i)z = 4+2i$, en donnant le résultat sous forme algébrique.
\end{exercice}
